\documentclass[11pt,a4paper]{article}
\usepackage[utf8]{inputenc}
\usepackage[T1]{fontenc}
\usepackage{amsmath,amssymb}
\usepackage{graphicx}
\usepackage{hyperref}
\usepackage{booktabs}
\usepackage{geometry}
\geometry{margin=2.5cm}

\title{Dynamic Object Motion Analysis (DOMA)\\Rapport de projet : classification de gestes à partir de pose et flot optique}
\author{Projet DOMA}
\date{\today}

\begin{document}
\maketitle

\begin{abstract}
Ce rapport décrit le pipeline \emph{Dynamic Object Motion Analysis} (DOMA) pour la reconnaissance de gestes de la main à partir de vidéos. Le système combine détection de la main (MediaPipe / YOLO), extraction de flot optique (Farnebäck / RAFT), et plusieurs modèles de classification séquence$\to$classe entraînés sur le jeu IPN Hand. Nous décrivons les méthodes (CNN-LSTM, Temporal Transformer, ST-GCN, ST-GCN-Opt), les formules de représentation des entrées (pose, métriques de flot), et nous comparons les performances sur l’ensemble de validation (exactitude, F1 macro/pondéré, perte). Les meilleurs résultats sont obtenus avec CNN-LSTM et Temporal Transformer sur les features unifiées pose+optflow.
\end{abstract}

\tableofcontents
\newpage

%==============================================================================
\section{Introduction}
%==============================================================================

La reconnaissance continue de gestes de la main (HGR) à partir de vidéos est utile pour les interfaces sans contact, la traduction de signes et le contrôle par gestes. Le projet DOMA vise à :

\begin{itemize}
  \item restreindre l’analyse du mouvement à une région d’intérêt (ROI) définie par la détection de main (ROI manuelle, MediaPipe Hands ou YOLO) ;
  \item calculer le flot optique dense dans cette ROI et en extraire des \emph{métriques} (vitesse, direction dominante, concentration) ;
  \item combiner ces métriques avec des descripteurs de \emph{pose} (position, vitesse, accélération du track 3D, optionnellement 21 landmarks) en un vecteur de features par frame ;
  \item entraîner des classificateurs séquence$\to$classe sur le benchmark IPN Hand et comparer plusieurs architectures.
\end{itemize}

Le rapport est structuré comme suit : Section~\ref{sec:data} présente le jeu de données IPN Hand ; Section~\ref{sec:pipeline} décrit l’extraction des features (pose et flot optique) et les formules associées ; Section~\ref{sec:methods} détaille les quatre modèles (CNN-LSTM, Temporal Transformer, ST-GCN, ST-GCN-Opt) ; Section~\ref{sec:experiments} donne le protocole expérimental et les tableaux de performances ; Section~\ref{sec:conclusion} conclut.

%==============================================================================
\section{Jeu de données IPN Hand}
\label{sec:data}
%==============================================================================

Le dataset \textbf{IPN Hand}~\cite{ipnhand} est un benchmark vidéo pour la reconnaissance continue de gestes de la main : plus de 4\,000 séquences de gestes, 800\,000 frames, 50 sujets, résolution 640$\times$480 à 30\,fps. Il comporte \textbf{14 classes} : une classe \emph{non-geste} (D0X) et 13 gestes (B0A, B0B, G01--G11) conçus pour l’interaction avec des écrans sans contact (pointing, click, throw, zoom, etc.).

Chaque sample est un segment vidéo annoté avec un label. Le pipeline DOMA produit, pour chaque sample, deux fichiers NPZ à partir des vidéos et annotations :

\begin{itemize}
  \item \textbf{Pose / cinématique} (\texttt{pose\_*.npz}) : position, vitesse et accélération du track 3D (\texttt{track\_pos\_xyz}, \texttt{track\_vel\_xyz}, \texttt{track\_acc\_xyz}), et optionnellement les 21 landmarks main (\texttt{landmarks\_xyz}) fournis par MediaPipe.
  \item \textbf{Flot optique} (\texttt{optflow\_*.npz}) : métriques temporelles par frame (vitesse moyenne, vitesse max, angle dominant, concentration de direction, seuil).
\end{itemize}

Un fichier \textbf{manifest} (\texttt{manifest.csv}) regroupe tous les samples avec les colonnes \texttt{sample\_id}, \texttt{split}, \texttt{label}, \texttt{pose\_npz}, \texttt{optflow\_npz}. Les modèles sont entraînés en utilisant soit un vecteur unifié pose+optflow par pas de temps, soit des représentations dérivées (skeleton+motion, track) pour les architectures à graphe ou multi-branches.

%==============================================================================
\section{Pipeline d’extraction des features}
\label{sec:pipeline}
%==============================================================================

\subsection{Flot optique et métriques}

Le flot optique dense fournit, pour chaque pixel, un vecteur de déplacement $\mathbf{f} = (u, v)$. À partir du champ de flot dans la ROI main, on calcule la \emph{magnitude} et l’\emph{angle} :
\begin{equation}
  m = \sqrt{u^2 + v^2}, \qquad \theta = \operatorname{atan2}(v, u) \in [-\pi, \pi].
\end{equation}
Un masque de mouvement est obtenu par seuillage de $m$ (seuil fixe, Otsu ou MAD). Après filtrage morphologique et sélection de la plus grande composante connexe, on soustrait éventuellement le mouvement global (médiane de $(u,v)$ en arrière-plan) pour limiter l’effet de la caméra.

Sur les pixels du masque, on définit les métriques suivantes, stockées dans \texttt{optflow\_*.npz} :
\begin{itemize}
  \item \textbf{Vitesse moyenne} et \textbf{vitesse max} : $\bar{m}$, $\max m$.
  \item \textbf{Direction dominante} (moyenne circulaire pondérée par la magnitude) :
  \begin{equation}
    C = \sum_i w_i \cos\theta_i, \quad S = \sum_i w_i \sin\theta_i, \quad \theta^* = \operatorname{atan2}(S, C), \quad R = \frac{\sqrt{C^2+S^2}}{\sum_i w_i}.
  \end{equation}
  $R \in [0,1]$ est la \emph{concentration} de direction ; $\theta^*$ est converti en degrés pour le stockage.
  \item \textbf{Seuil} utilisé pour le masque (Otsu, MAD ou fixe).
\end{itemize}

\subsection{Représentation pose}

À chaque frame $t$, la pose est représentée par :
\begin{itemize}
  \item \textbf{Track 3D} : position $\mathbf{p}_t$, vitesse $\mathbf{v}_t$, accélération $\mathbf{a}_t$ (concaténés en un vecteur de dimension 9 par frame pour le modèle ST-GCN-Opt).
  \item \textbf{Skeleton + motion} (pour ST-GCN) : soit les 21 landmarks $\mathbf{L}_t \in \mathbb{R}^{21 \times 3}$ (transposés en tenseur $3 \times T \times 21$), avec motion $\mathbf{M}_t = \mathbf{L}_t - \mathbf{L}_{t-1}$ (et zéro à $t=0$), soit, en l’absence de landmarks, le track seul (position en skeleton, vitesse en motion), avec un seul nœud ($n=1$).
\end{itemize}

Les entrées des modèles sont donc toutes dérivées des deux sources \textbf{pose} et \textbf{optflow} ; selon l’architecture, on utilise soit un vecteur concaténé par frame (Temporal Transformer, CNN-LSTM), soit skeleton+motion et/ou track+optflow (ST-GCN, ST-GCN-Opt).

%==============================================================================
\section{Méthodes (modèles)}
\label{sec:methods}
%==============================================================================

On note $\mathcal{X} = (\mathbf{x}_1, \ldots, \mathbf{x}_T)$ la séquence d’entrée (chaque $\mathbf{x}_t$ est un vecteur de features ou un tenseur selon le modèle), $\ell$ les longueurs réelles par sample, et $K$ le nombre de classes. L’objectif est d’estimer une distribution sur les classes $\hat{y} = \operatorname{argmax}_k \mathbf{z}_k$ à partir des logits $\mathbf{z} \in \mathbb{R}^K$. Tous les modèles sont entraînés par minimisation de la perte de cross-entropy :
\begin{equation}
  \mathcal{L} = -\frac{1}{N}\sum_{n=1}^{N} \log \frac{\exp(z_{n,y_n})}{\sum_{k=1}^{K} \exp(z_{n,k})}.
\end{equation}

\subsection{CNN-LSTM}

L’entrée est une séquence de vecteurs $\mathbf{x}_t \in \mathbb{R}^F$ (pose + optflow unifiés, $F=78$ par défaut). Le modèle comporte :
\begin{enumerate}
  \item \textbf{Conv1D temporelle} : couches de convolution 1D le long du temps (noyau $k$, canaux $C$), BatchNorm, GELU, Dropout, produisant $\mathbf{H}^{(1)} \in \mathbb{R}^{B \times T \times C}$.
  \item \textbf{LSTM bidirectionnelle} : $\overrightarrow{\mathbf{h}}_t$, $\overleftarrow{\mathbf{h}}_t$ ; sortie concaténée par pas de temps.
  \item \textbf{Pooling temporel masqué} : moyenne sur les pas valides (non padding) en fonction des longueurs $\ell$ :
  \begin{equation}
    \mathbf{h} = \frac{\sum_{t=1}^{T} \mathbf{m}_t \cdot \mathbf{h}_t}{\sum_{t} m_t}, \quad m_t = \mathbb{1}[t \le \ell].
  \end{equation}
  \item \textbf{Couche de sortie} : LayerNorm, Dropout, Linear($C_{\mathrm{lstm}}$, $K$) $\to$ logits $\mathbf{z}$.
\end{enumerate}

\subsection{Temporal Transformer}

L’entrée est la même séquence $\mathbf{x}_t \in \mathbb{R}^F$. Le modèle :
\begin{enumerate}
  \item \textbf{Normalisation et projection} : LayerNorm($F$), Linear($F \to d$), LayerNorm($d$), avec $d$ la dimension du modèle.
  \item \textbf{Encodage positionnel sinusoïdal} :
  \begin{equation}
    \mathrm{PE}(t, 2i) = \sin(t / 10000^{2i/d}), \quad \mathrm{PE}(t, 2i+1) = \cos(t / 10000^{2i/d}).
  \end{equation}
  \item \textbf{Transformer Encoder} : couches d’attention multi-têtes et FFN (GELU, norm-first), avec masque de padding pour ignorer les pas au-delà de $\ell$.
  \item \textbf{Pooling temporel masqué} : moyenne des sorties du Transformer sur les positions valides, puis LayerNorm et classifieur linéaire $d \to K$.
\end{enumerate}

\subsection{ST-GCN (Spatial-Temporal Graph Convolutional Network)}

L’entrée est un tenseur $\mathbf{X} \in \mathbb{R}^{B \times 6 \times T \times n}$ (skeleton et motion concaténés : 3 canaux skeleton + 3 canaux motion, $T$ pas de temps, $n$ nœuds, $n=21$ avec landmarks). Le graphe spatial est défini par une matrice d’adjacence $\mathbf{A} \in \mathbb{R}^{n \times n}$ (ici $\mathbf{A} = \mathbf{I}$ pour un graphe sans arêtes explicites, ou une matrice de voisinage apprise).

\paragraph{Bloc spatio-temporel (ST-Block).}
\begin{enumerate}
  \item \textbf{Convolution temporelle} (noyau $k_t$) avec activation GLU : 
  $\mathbf{X}' = (\mathbf{W}_t \ast \mathbf{X}_{:,:,1:c_1,:} + \mathbf{b}) \odot \sigma(\mathbf{W}_t' \ast \mathbf{X})$ (partie gate).
  \item \textbf{Convolution spatiale (graph)} : pour un canal $c$, 
  $\mathbf{Y}_{b,c,t,n} = \sum_{m} \mathbf{A}_{n,m} \mathbf{X}_{b,c,t,m}$ puis combinaison linéaire avec noyau spatial $\boldsymbol{\Theta}$ et biais, puis ReLU et connexion résiduelle.
  \item \textbf{Deuxième convolution temporelle} (ReLU), LayerNorm sur les dimensions ($n$, $c$), Dropout.
\end{enumerate}

Deux ST-Blocks sont empilés, suivis d’un \textbf{pooling global} (moyenne sur temps et nœuds) et d’un classifieur linéaire vers $K$ classes.

\subsection{ST-GCN-Opt (multi-branches)}

Ce modèle fusionne trois branches à partir des mêmes données pose + optflow :
\begin{enumerate}
  \item \textbf{Branche ST-GCN} : comme ci-dessus, sur skeleton+motion $\to$ vecteur $\mathbf{h}_{\mathrm{stgcn}}$.
  \item \textbf{Branche track} : séquence track (pos, vel, acc) de dimension 9 par frame ; CNN 1D temporelle (Conv1D + BN + ReLU + Dropout) + pooling global $\to$ $\mathbf{h}_{\mathrm{track}}$.
  \item \textbf{Branche optflow} : séquence des 6 métriques de flot par frame ; même structure CNN 1D $\to$ $\mathbf{h}_{\mathrm{opt}}$.
\end{enumerate}

La fusion peut être \emph{concaténation}, \emph{pondérée} (poids appris + softmax), \emph{attention} (projection commune + multi-head attention) ou \emph{gated}. Les vecteurs fusionnés sont passés dans un MLP (Linear$\to$BN$\to$ReLU$\to$Dropout, répété) puis Linear $\to$ logits $K$.

%==============================================================================
\section{Expériences et résultats}
\label{sec:experiments}
%==============================================================================

\subsection{Protocole}

\begin{itemize}
  \item \textbf{Données} : manifest IPN Hand, split train/val/test ; les métriques rapportées sont sur l’ensemble de \textbf{validation} (best epoch selon l’exactitude sur la validation).
  \item \textbf{Entraînement} : 20 epochs, AdamW (lr $5\times10^{-4}$, weight decay 0.01), CosineAnnealingLR, batch size 32, gradient clipping (max norm 1). Une seule graine utilisée pour les runs rapportés.
  \item \textbf{Métriques} : exactitude (accuracy), précision/rappel/F1 en macro et pondéré (weighted).
\end{itemize}

\subsection{Comparaison des modèles}

Les tableaux suivants résument les \textbf{meilleures métriques} (best epoch) pour chaque modèle sur la validation. Les runs proviennent des dossiers \texttt{models/<model>\_<date>\_<heure>/} (metrics.json).

\begin{table}[htbp]
  \centering
  \caption{Comparaison des performances (validation) — métriques globales.}
  \label{tab:main}
  \begin{tabular}{lcccccc}
    \toprule
    Modèle & Accuracy & P (macro) & R (macro) & F1 (macro) & F1 (weighted) & Val. loss \\
    \midrule
    CNN-LSTM           & 0.883 & 0.817 & 0.836 & 0.823 & 0.884 & 0.595 \\
    Temporal Transformer & 0.850 & 0.742 & 0.764 & 0.748 & 0.850 & 0.610 \\
    ST-GCN-Opt         & 0.665 & 0.465 & 0.300 & 0.311 & 0.597 & 1.365 \\
    ST-GCN             & 0.538 & 0.246 & 0.226 & 0.196 & 0.458 & 1.549 \\
    \bottomrule
  \end{tabular}
\end{table}

\begin{table}[htbp]
  \centering
  \caption{Perte d’entraînement et epoch du meilleur modèle (validation).}
  \label{tab:train}
  \begin{tabular}{lcc}
    \toprule
    Modèle & Train loss (best) & Best epoch \\
    \midrule
    CNN-LSTM            & 0.0629 & 20 \\
    Temporal Transformer & 0.3464 & 18 \\
    ST-GCN-Opt          & 0.8283 & 17 \\
    ST-GCN              & 1.6031 & 20 \\
    \bottomrule
  \end{tabular}
\end{table}

\subsection{Analyse}

\begin{itemize}
  \item \textbf{CNN-LSTM} et \textbf{Temporal Transformer} obtiennent les meilleures performances (exactitude $\approx$ 85--88\%, F1 macro $\approx$ 0.75--0.82), en utilisant directement le vecteur unifié pose+optflow par frame. La représentation séquentielle standard convient bien au jeu IPN Hand avec les features extraites.
  \item \textbf{ST-GCN-Opt} est en retrait (accuracy $\approx$ 66.5\%, F1 macro $\approx$ 0.31), avec une forte variance sur les classes de gestes (bon rappel sur D0X/B0A/B0B, rappel faible sur plusieurs G0x). La fusion multi-branches et le format skeleton/track/optflow pourraient nécessiter un réglage (learning rate, régularisation, ou plus de données) pour converger mieux.
  \item \textbf{ST-GCN} seul atteint seulement $\approx$ 53.8\% d’exactitude (F1 macro $\approx$ 0.20), avec de nombreux rappels à zéro sur les gestes G04--G06, G10, G11. L’utilisation du graphe skeleton+motion sans les branches track/optflow et avec un graphe trivial ($\mathbf{A}=\mathbf{I}$) limite probablement la capacité du modèle sur ce dataset.
\end{itemize}

En résumé, pour ce pipeline et ce split, les modèles basés sur la séquence de features unifiées (CNN-LSTM, Temporal Transformer) surpassent les modèles à graphe spatio-temporel (ST-GCN, ST-GCN-Opt) dans la configuration actuelle.

%==============================================================================
\section{Conclusion}
\label{sec:conclusion}
%==============================================================================

Le projet DOMA fournit un pipeline complet pour la reconnaissance de gestes de la main à partir de vidéos : détection de la main (MediaPipe/YOLO), flot optique (Farnebäck/RAFT), extraction de métriques (magnitude, direction dominante, concentration) et de pose (track 3D, landmarks), puis classification par plusieurs modèles (CNN-LSTM, Temporal Transformer, ST-GCN, ST-GCN-Opt). Les formules de représentation (moyenne circulaire, pooling masqué, blocs ST-GCN) ont été rappelées et les performances sur IPN Hand comparées en tableaux. Les meilleurs résultats sont obtenus avec CNN-LSTM et Temporal Transformer sur les features pose+optflow unifiées ; les architectures à graphe pourraient être améliorées par un graphe spatial plus riche et un réglage dédié.

%==============================================================================
\section*{Références}
%==============================================================================

\begin{enumerate}
  \item \label{ipnhand} G. Benitez-Garcia et al., ``IPN Hand: A Video Dataset and Benchmark for Real-Time Continuous Hand Gesture Recognition,'' in \textit{ICPR}, 2021.
  \item G. Farnebäck, ``Two-frame motion estimation based on polynomial expansion,'' in \textit{Scandinavian Conference on Image Analysis}, 2003.
  \item Y. Yan, S. Xiong, et D. Lin, ``Spatial temporal graph convolutional networks for skeleton-based action recognition,'' in \textit{AAAI}, 2018.
  \item A. Vaswani et al., ``Attention is all you need,'' in \textit{NeurIPS}, 2017.
\end{enumerate}

\end{document}
